\section[Une institution juridique]{Une institution au cœur de l’écosystème de la recherche juridique}

Depuis 1979, la bibliothèque Cujas, devenu bibliothèque interuniversitaire, est gérée par une co-tutellle. L'Université Paris I Panthéon-Sorbonne s'occupe de l'administratif et l'Université Paris II Panthéon-Assas régi l'établissement de cette institution. Le public consultant ses collections se constitue \enquote{d'étudiants de L2 aux enseignants-chercheurs et professionnels du droit}. Cela représente 16 500 usagers inscrits à la bibliothèque. 

La bibliothèque Cujas s'intègre dans différents réseaux documentaires nationaux lui permettant d'acquérir le statut de pôle d'excellence. Elle est également connue pour être un pôle associé de la Bibliothèque nationale de France (BNF), ensemble ils copilotent \enquote{le programme national de numérisation concertée en sciences juridiques}. La bibliothèque Cujas est partenaire de longue date de la BnF. En s'associant à la BNF sur le programme de numérisation, la bibliothèque Cujas espère donner une meilleure visibilité de la richesse de ses collections en ligne. En effet, ce partenariat ancien et solide entre les deux institutions vise à \enquote{créer un réseau documentaire thématique, rassemblant des institutions de tout statut et de toute taille qui mènent des projets en matière de numérisation du patrimoine juridique}. 

Lionel Maurel, conservateur à la BNF, se charge de la phase prospection du projet afin d'identifier tous les \enquote{gisements documentaires en droit et des bibliothèques numériques existantes}. Ensuite, cette cartographie se transforme en \enquote{deux appels à initiatives pour la numérisation de corpus imrpimés, articulés autour des quatre sources traditionnelles du droit : sources législatives et réglementaires, jurisprudence, doctrine, sources du droit local}. Grâce au soutien de la BNF, les collections numérisées pour ces appels à projet permettent une excellente visibilité sur la bibliothèque numérique de la BNF, Gallica. 

La bibliothèque Cujas s'intègre dans d'autres réseaux nationaux dans le réseau CollEx. Cujas est officiellement reconnue comme \enquote{bibliothèque délégataire CollEx en droit et sciences juridiques}, l'établissant comme l'une des têtes du réseau nationale de l'Enseignement Supérieur et de la Recherche (ESR). Pour renforcer son interconnexion avec les autres bibliothèques, Cujas a entrepris de signaler une partie de ses collections dans Calames pour ses fonds archivistiques comme le fonds ONU. Elle participe activement au signalement des périodiques via Le Doctrinal et en participant au projet HOPPE-Droit. Ces différents projets renforcent la position centrale de la bibliothèque Cujas dans le domaine juridique.

A cette époque, le contexte budgétaire était appauvrie, la numérisation des collections de la bibliothèque Cujas à travers l'utilisation des chaînes de numérisation de la BNF. En profitant de la chaîne de numériation de la BNF, celle-ci intègre les collections des partenaires à l'intérieur de Gallica, leur permettant de profiter de la visibilité de Gallica pour mettre en valeur les collections provenant des autres institutions, en écrivant l'origine du document. Cela implique un meilleur signalement des collections numérisées via le moissonnage OAI-PMH, condition requise pour permettre le rapprochement virtuel des collections numérisées.